\documentclass[12pt,a4paper]{article}
\usepackage[T1]{fontenc}
\usepackage[utf8]{inputenc}
\usepackage[polish]{babel}
\usepackage[a4paper,margin=2.5cm]{geometry}
\usepackage{amsmath,amssymb}
\usepackage{listings}

\linespread{0.96} % lekko ciaśniej (bez psucia czytelności)

\title{\textbf{Sprawozdanie 1}\\Wybrane Zagadnienia Algebry}
\author{Jakub Kowal\\279706}
\date{}

\begin{document}
\maketitle

Użyte parametry do testów:
\[
a=2,\ b=7,\ c=9,\ d=7,\ e=0,\ f=6.
\]

\section*{1. Pierscień liczb Gaussa $\mathbb{Z}[i]$}

\subsection*{Reprezentacja}
Liczba Gaussa jest reprezentowana jako para liczb całkowitych \texttt{Pair (a, b)},
czyli element $a+bi$.
W pliku kodzie:
\begin{itemize}
\item \texttt{Pair.first} to czesc rzeczywista,
\item \texttt{Pair.second} to czesc urojona.
\end{itemize}

\subsection*{Norma}
Norma jest liczona jako
\[
N(a+bi)=a^2+b^2.
\]

\subsection*{Dzielenie z resztą}
Dzielenie jest realizowane przez przybliżenie ilorazu zespolonego i sprawdzenie
możliwych zaokrągleń (floor/ceil) dla części rzeczywistej i urojonej.
Reszta jest poprawna, gdy spełnia $N(r) < N(b)$.

Fragment kodu:
\begin{verbatim}
function Dziel_wszystkie(a::Pair{<:Integer,<:Integer},b::Pair{<:Integer,<:Integer})
	denominator = float(Norma(b))
	real = (a.first*b.first + a.second*b.second) / denominator
	imag = (a.second*b.first - a.first*b.second) / denominator

	potential_reals = Set([floor(Int, real), ceil(Int, real)])
	potential_imags = Set([floor(Int, imag), ceil(Int, imag)])

	wyniki = Vector{Tuple{Pair{Int,Int},Pair{Int,Int}}}()
	for r in potential_reals
		for i in potential_imags
			q = Pair(r, i)
			rem = Odejmij(a, Mnoz(b, q))
			if Norma(rem) < Norma(b)
				push!(wyniki, (q, rem))
			end
		end
	end
	return wyniki
end
\end{verbatim}

\subsection*{NWD i NWW}
NWD jest liczony algorytmem Euklidesa z dzieleniem w $\mathbb{Z}[i]$,
a NWW ze wzoru
\[
\mathrm{NWW}(a,b)=\frac{ab}{\mathrm{NWD}(a,b)}.
\]

Fragment kodu NWD:
\begin{verbatim}
function NWD(a::Pair{<:Integer,<:Integer},b::Pair{<:Integer,<:Integer})
	while !(b.first == 0 && b.second == 0)
		_,r=Dziel(a,b)
		a=b
		b=r
	end
	return a
end
\end{verbatim}

\subsection*{Wyniki wywolan z TestyZ1.jl}
Uzyte parametry:
\[
a=2,\ b=7,\ c=9,\ d=7,\ e=0,\ f=6.
\]

\begin{itemize}
\item Norma:
\[
N(a-bi)=N(2-7i)=2^2+7^2=53.
\]

\item Dzielenie $(c+a)+(d+b)i$ przez $e+fi$:
\[
\frac{11+14i}{6i}.
\]
Wyniki z testu (wszystkie znalezione poprawne pary $(q,r)$):
\begin{enumerate}
    \item q=$2-i$, r=$5+2i$
    \item q=$2-2i$, r=$-1+2i$
    \item q=$3-3i$, r=$-1-4i$
\end{enumerate}

\item Dlaczego nie ma ich wiecej:
z metody floor/ceil jest najwyzej 4 kandydatów na iloraz,
ale tylko ci, ktorzy dają $N(r)<N(b)$, są akceptowani.
Tutaj 3 kandydatów spełnia warunek.

\item NWD i NWW dla trojki $a+bi, c+di, e+di$:
\begin{verbatim}
NWD([a+bi, c+di, e+di]) = 1
NWW([a+bi, c+di, e+di]) = -217 + 539i
\end{verbatim}
Żeby dostać wszystkie wersje w $\mathbb{Z}[i]$ trzeba wymnożyć te wyniki o jednostki $\{1,-1,i,-i\}$.

\item Wywolania listowe:
\begin{verbatim}
NWD([a+bi]) = 2 + 7i
NWD([])     = 0 + 0i
\end{verbatim}
W kodzie obsluga pustej listy dla NWD/NWW zwraca $0+0i$.
Dla listy 1-elementowej funkcje zwracaja ten element.
\end{itemize}

\section*{2. Pierścień wielomianów $\mathbb{R}[x]$}

\subsection*{Reprezentacja i rozszerzalność}
Wielomian jest reprezentowany przez
\begin{verbatim}
struct Pierscien{T}
	coeffs::Dict{Vector{Int}, T}
end
\end{verbatim}
Klucz to wektor wykładników, a wartość to współczynnik.
W $\mathbb{R}[x]$ klucz ma postać jednowymiarową, np. \texttt{[3]} dla $x^3$.
Ta sama struktura łatwo przechodzi na $\kappa[x_1,\dots,x_n]$,
bo wtedy klucz to po prostu wektor długości $n$.

\subsection*{Norma, dzielenie, NWD, NWW, rozszerzony NWD}
\begin{itemize}
\item Norma wielomianu: bierzemy maksymalny wykładnik przy niezerowym współczynniku (stopień wielomianu), a dla wielomianu zerowego zwracamy $-1$.
\item Dzielenie z resztą: w każdej iteracji bierzemy iloraz wyrazów wiodących, tworzymy z niego kolejny składnik ilorazu i odejmujemy odpowiednią wielokrotność dzielnika od bieżącej reszty, aż stopień reszty będzie mniejszy od stopnia dzielnika.
\item NWD: stosujemy algorytm Euklidesa, czyli powtarzamy przejście $(a,b) \mapsto (b, a\bmod b)$ do uzyskania zera.
\item NWW: liczymy iloczyn wielomianów i dzielimy go przez NWD.
\item Rozszerzony NWD: równolegle z algorytmem Euklidesa aktualizujemy współczynniki i zwracamy trójkę $(d,s,t)$ taką, że
\[
d = s\,a + t\,b.
\]
\end{itemize}

\subsection*{Wyniki wywołań z TestyZ2.jl}

\paragraph{a) Norma i dzielenie}
Norma (stopień) dla naszych danych wynosi:
\[
Norma(9x^2+7)=2
\]

Wynik dzielenia z resztą przez $x+1$ (dokładnie taki dzielnik jest w pliku testowym):
\[
\frac{9x^2+7}{x + 1}=9x-9 \text{z resztą:} 16
\]
\begin{verbatim}
(Pierscien{Float64}(Dict([1] => 9.0, [0] => -9.0)),
 Pierscien{Float64}(Dict([0] => 16.0)))
\end{verbatim}
czyli
\[
9x^2+7 = (x+1)(9x-9) + 16.
\]

\paragraph{b) Rozszerzony NWD dla pary $v,w$ oraz wyznaczenie $g$}
W teście:
\[
v(x)=2x^3+7x^2+9x+7,
\quad
w(x)=7x^3+6x.
\]

Z rozszerzonego NWD:
\begin{verbatim}
(Pierscien{Float64}(Dict([0] => 1.0)),
 Pierscien{Float64}(Dict([1] => -0.15668693962003885, [0] => 0.14285714285714288,
						 [2] => 0.1415762743745689)),
 Pierscien{Float64}(Dict([1] => -0.09680857734027208, [0] => -0.03148428472900224,
						 [2] => -0.04045036410701969)))
\end{verbatim}
Czyli otrzymujemy, że NWD(v,w) jest 1, czyli że są względnie pierwsze.

Skoro są względnie pierwsze to szukamy stałej $g \in \mathbb{R}$ takiej, że
\[
Norma(NWD(v, w+g)) > 0,
\]
i $v$ oraz $w+g$ bedą miały wspólny pierwiastek x0. Jeśli x0 jest pierwiastkiem $v$, to wystarczy dobrać
\[
g = -w(x0).
\]
W implementacji najpierw przybliżamy pierwiastek x0 wielomianu $v$ metodą stycznych (metoda Newtona z Listy 3 na labolatoria z Obliczeń Naukowych), a potem liczymy $g=-w(x0)$.

Wyznaczone zostało
\[
g = 84.48215863265216,
\]
więc nowe w to:
\[
w(x)+g = 7x^3 + 6x + 84.48215863265216.
\]

Następnie:
\begin{verbatim}
NWD(v, w+g) = Pierscien{Float64}(Dict([1] => 1.0, [0] => 2.1693702699113073))
\end{verbatim}
czyli
\[
\mathrm{NWD}(v,w+g)=x+2.1693702699113073,
\]
więc rzeczywiście $1$ nie jest NWD.

Odpowiadające NWW:
\begin{verbatim}
Pierscien{Float64}(Dict([5] => 1.0, [4] => 1.330629730088693,
						[3] => 2.4705142804283398, [2] => 13.20941957331203,
						[1] => 17.442100068341315, [0] => 19.47158578791304))
\end{verbatim}
czyli
\begin{align*}
\mathrm{NWW}(v,w+g) = {} & x^5 + 1.330629730088693x^4 \\
& + 2.4705142804283398x^3 \\
& + 13.20941957331203x^2 \\
& + 17.442100068341315x \\
& + 19.47158578791304
\end{align*}

\section*{3. Porządki produktowe w $\mathbb{N}^{n}$}
Sprawdzane pary w $\mathbb{N}^{2}$:
\begin{itemize}
    \item (2,7)
    \item (9,7)
    \item (0,6)
\end{itemize}
Wyniki dla par:
\[
\begin{array}{c}
\text{Porównanie} \\\hline
(2,7)\le(9,7) \\
(2,7)\not \le(0,6)\\
(9,7)\not \le(2,7)\\
(9,7)\not \le(0,6) \\
(0,6)\le(2,7) \\
(0,6)\le(9,7)
\end{array}
\]

Sprawdzane trójki w $\mathbb{N}^{3}$:
\begin{itemize}
    \item (2,9,0)
    \item (7,7,6)
\end{itemize}
Wyniki dla trójek:

\[
\begin{array}{c}
    \text{Porównanie}\\\hline
(2,9,0)\not \le(7,7,6)\\
(7,7,6)\not \le(2,9,0)
\end{array}
\]

Punkty minimalne znalezione w A = \{(x,y)$\in \mathbb{N}^{2}$: $(x-2)^{2}+(y-7)^{2}$<5\}:
\begin{itemize}
    \item (0,6)
    \item (1,5)
\end{itemize}

Lista punktów znalezionych w B = \{($x_{1},x_{2},x_{3},x_{4}$)$\in \mathbb{N}^{4}$: $(x_{1}-9)^{2} + (x_{2}-7)^{2} + (x_{3}-0)^{2} + (x_{4} - 6)^{2} > 224 $\}
{\scriptsize
\begin{lstlisting}
[0,0,0,16], [0,0,4,15], [0,0,6,14], [0,0,7,13], [0,0,8,0],
[0,15,0,15], [0,15,4,14], [0,15,6,13], [0,15,7,0],
[0,16,0,14], [0,16,4,13], [0,16,6,0], [0,17,0,13],
[0,17,3,0], [19,0,0,15], [19,0,4,14], [19,0,6,13],
[19,0,7,0], [19,15,0,14], [19,15,4,13], [19,15,5,0],
[19,16,0,13], [19,16,3,0], [19,17,0,0]
\end{lstlisting}
}

\section*{Wykorzystanie Sztucznej Inteligencji}
Wykorzystałem sztuczną inteligencję w ramach wsparcia redakcji opisów implementacji oraz usprawnienia przeszukiwania dokumentacji języka Julia
\end{document}
